% =========================================================
% main.tex — Polar Semirings
% Structure fixed by final-rules.md (5 sections)
% =========================================================
\documentclass[11pt,a4paper]{article}

% =========================================================
% head-tools.tex — packages, environments, macros
% =========================================================

% ---------- Block 1: encoding & fonts ----------
% On full TeX installations lmodern + T1 are used; on minimal ones we fall
% back to the default OT1/CM setup to avoid bitmap-font generation.
\IfFileExists{lmodern.sty}
  {\usepackage[T1]{fontenc}\usepackage{lmodern}}
  {}
\usepackage[utf8]{inputenc}
\IfFileExists{lmodern.sty}{\usepackage{microtype}}{\usepackage[expansion=false]{microtype}}

% ---------- Block 2: mathematics ----------
\usepackage{amsmath}
\usepackage{amssymb}
\usepackage{amsthm}
\usepackage{mathtools}

% ---------- Block 3: tables & layout ----------
\usepackage{booktabs}
\usepackage{array}
\usepackage[margin=2.7cm]{geometry}

% ---------- Block 4: bibliography ----------
\usepackage[numbers,sort&compress]{natbib}

% ---------- Block 5: hyperref & cleveref (load last) ----------
\usepackage[colorlinks=true,linkcolor=blue!60!black,citecolor=green!45!black,urlcolor=blue!60!black]{hyperref}
\usepackage[capitalise,noabbrev]{cleveref}

% ---------- Block 6: theorem environments ----------
\theoremstyle{plain}
\newtheorem{theorem}{Theorem}[section]
\newtheorem{proposition}[theorem]{Proposition}
\newtheorem{lemma}[theorem]{Lemma}
\newtheorem{corollary}[theorem]{Corollary}

\theoremstyle{definition}
\newtheorem{definition}[theorem]{Definition}
\newtheorem{example}[theorem]{Example}

\theoremstyle{remark}
\newtheorem{remark}[theorem]{Remark}

\crefname{axiom}{axiom}{axioms}

% ---------- Block 7: Polar Semiring notation ----------
% Primitive: \vee, \otimes, ^{*}, e_\vee, e_\otimes
% Derived:   \wedge, \oplus, e_\wedge, e_\oplus
\newcommand{\ev}{e_{\vee}}
\newcommand{\ew}{e_{\wedge}}
\newcommand{\eo}{e_{\otimes}}
\newcommand{\ep}{e_{\oplus}}
\newcommand{\PS}{\mathrm{PS}}
\newcommand{\BGA}{\mathrm{BGA}}
\newcommand{\dstar}[1]{{#1}^{*}}          % polarity applied to #1
\newcommand{\leqm}{\leq^{*}}              % meet order
% Residual (BGA-native; permitted ONLY in Section 4, cf. final-rules §4)
\newcommand{\res}{\mathbin{\to}}

% ---------- Block 8: author-check notes ----------
% \noteAI{...} marks items the author must verify before submission.
% Set \noteAIfalse to suppress all notes in the final version.
\newif\ifnoteAI
\noteAItrue
\usepackage{xcolor}
\ifnoteAI
  \newcommand{\noteAI}[1]{\textcolor{red!70!black}{\footnotesize [\textsf{AI-check}: #1]}}
\else
  \newcommand{\noteAI}[1]{}
\fi

% ---------- Block 9: architecture diagram ----------
\usepackage{tikz}
\usetikzlibrary{positioning,fit,backgrounds,arrows.meta,calc}


\title{Polar Semirings}
\author{%
    Thu-Le Tran\\
    Can Tho University \\
    ttle@ctu.edu.vn
}
\date{\today}

\begin{document}

\maketitle

\begin{abstract}
We introduce \emph{Polar Semirings}, a class of algebraic structures in which
polarity---an order-reversing involution---is taken as a primitive operation
rather than derived from residuation. A Polar Semiring is axiomatized by
eleven equational and quasi-equational axioms over the primitive signature
$(X,\vee,\otimes,{}^{*},e_\vee,e_\otimes)$; the meet, the plus operation, and
their units are obtained by conjugation under the polarity. We show that the
axiom system is independent, exhibit canonical models, and prove a
Fundamental Structure Theorem: every Polar Semiring canonically induces a dual
Polar Semiring, with the polarity acting as an isomorphism and the two induced
orders coinciding. We then position the class within the landscape of
involutive algebraic structures: every bounded Girard algebra induces a Polar
Semiring, whereas there exists a Polar Semiring on the unit interval that
cannot be extended to a bounded Girard algebra. Consequently, bounded Girard
algebras form a proper subclass, and Polar Semirings constitute a genuine
residuation-free generalization.
\end{abstract}

\begin{center}
    \textbf{Keywords:} polarity; involution; semiring; Girard algebra; residuation
\end{center}

% \noteAI{Add MSC2020 codes (suggest: 06B05, 06D30, 16Y60, 03G10)}

% =========================================================
% doc1-intro.tex — Section 1: Introduction
% Source: d0 §Sec1; d1 §1 (motivation, contributions, related work)
% =========================================================

\section{Introduction}\label{sec:intro}

Order-reversing involutions arise in two largely separate traditions. In the
theory of residuated structures, of which bounded Girard algebras are the
representative example, a structure carries a negation $\dstar{(\cdot)}$
obtained from the residual by fixing a distinguished element $\ep$ and setting
$\dstar{a} := a \res \ep$ \citep{Rosenthal1990,Agliano2025}. This negation has
two properties: it is an involution, $\dstar{(\dstar{a})} = a$, and it is
antitone with respect to the order, $a \leq b$ implying
$\dstar{b} \leq \dstar{a}$. In this tradition the residual is the source of
both properties: the involution and its order-reversal are consequences of the
residuation law, not primitive assumptions.

In algebra, by contrast, a semiring with involution is equipped with a negation
$\dstar{(\cdot)}$ as a primitive operation satisfying only
$\dstar{(\dstar{a})} = a$, that is, an anti-automorphism of the semiring
\citep{Golan1999,dolinka2000minimal}. Here the negation is primitive rather than
manufactured, but it is not antitone; indeed the structure carries no order at
all with respect to which antitonicity could be stated.

These two traditions leave a natural question open: is there a structure in
which the negation is a primitive operation possessing both properties at
once---involution and antitonicity---without being generated from a residual?

This paper answers the question affirmatively. We retain the involution axiom
of the algebraic tradition and replace the residual by a single new axiom
relating the negation to the join. On this basis the order is not primitive but
derived: it is recovered from the join, and antitonicity is not assumed but
follows. Concretely, we introduce \emph{Polar Semirings}: idempotent
commutative semirings equipped with a primitive negation subject to two axioms,
the involution axiom and a new polarity--join compatibility axiom. On this
structure (i)~the order and the antitonicity of the negation are generated
rather than assumed; (ii)~no residual is used; and (iii)~in contrast to a
semiring with involution, the negation is not an anti-automorphism of a single
structure but an isomorphism onto a new dual idempotent commutative semiring.

This is not the only work to obtain a primitive, order-compatible polarity.
Dual pairs of semirings, developed to study fuzzy sets, provide such a polarity
as well \citep{movckovr2025unification}; that system, however, imposes
additional distributivity axioms tailored to its intended application. Our
concern is not a particular application but the axiomatic question of how little
is required: we add only the two axioms governing the negation on top of a
standard semiring, isolating the polarity and discarding the distributivity
constraints entirely. This isolation makes it possible (i)~to separate the
contribution of the polarity itself from any distributivity assumption; (ii)~to
identify the minimal axiomatic cost of a primitive, order-compatible polarity;
and (iii)~to obtain a structure that does not presuppose completeness or the
extra distributivity laws of dual pairs of semirings.

\paragraph{Contributions.}
The paper makes three contributions.
\begin{enumerate}
  \item \emph{Definition and minimality} (\Cref{sec:ps}). We axiomatize Polar
  Semirings by eleven axioms, exhibit canonical models, and prove that the
  axiom system is independent: no axiom is a consequence of the remaining
  ten.
  \item \emph{Architecture} (\Cref{sec:arch}). We show that the axioms
  decompose into three layers: an algebraic layer (the polarity-free
  semiring reduct), a reflection layer, in which the involution axiom alone
  generates the dual operations and all De Morgan laws at no axiomatic cost,
  and an integration layer, in which the single polarity axiom fuses the
  reflected order with the original one. The decomposition culminates in the
  Fundamental Structure Theorem: every Polar Semiring canonically induces a
  dual Polar Semiring, the polarity is an isomorphism onto the dual, and the
  two induced orders coincide.
  \item \emph{Positioning} (\Cref{sec:positioning}). We locate Polar
  Semirings relative to bounded Girard algebras. Every bounded Girard
  algebra, upon forgetting its residual, satisfies all eleven axioms; and
  there exists a Polar Semiring on the unit interval that cannot be extended
  to a bounded Girard algebra. Hence bounded Girard algebras form a proper
  subclass, and Polar Semirings are a strict, residuation-free
  generalization.
\end{enumerate}

\paragraph{Related structures.}
The order-theoretic reduct of a Polar Semiring is a bounded lattice with an
antitone involution satisfying the De Morgan laws, that is, a bounded
i-lattice in the sense of \citet{Kalman1958}; distributivity is not assumed,
so De Morgan algebras \citep{Rasiowa1974} and Boolean algebras arise from
Polar Semirings only after adding further axioms. On the multiplicative side,
semirings with involution \citep{Golan1999} retain the involution but drop
the order-compatibility axiom entirely. A brief orientation table comparing
these classes precedes the detailed comparison with bounded Girard algebras
in \Cref{sec:positioning}. Throughout the paper we maintain a strict
terminological distinction, elaborated in \Cref{sec:arch}: a structure is
called \emph{involutive} when it satisfies only the equation
$\dstar{(\dstar{a})}=a$, and \emph{polar} when the involution is additionally
compatible with the induced order, in Birkhoff's classical sense
\citep{Birkhoff1940}.

\paragraph{Outline.}
\Cref{sec:ps} defines Polar Semirings, presents basic models, and proves
independence of the axioms. \Cref{sec:arch} develops the three-layer
architectural analysis and proves the Fundamental Structure Theorem.
\Cref{sec:positioning} carries out the comparison with bounded Girard
algebras, proving both the inclusion and its strictness.
\Cref{sec:conclusion} summarizes and discusses future directions.

% =========================================================
% doc2-polar-semiring.tex — Section 2: Primitive Polarity on Semirings
% =========================================================

\section{Primitive Polarity on Semirings}\label{sec:ps}

The objective of this section is to establish primitive polarity as an
independent algebraic primitive that can be adjoined to an idempotent
commutative semiring without appeal to residuals or dual operations. We
first fix the underlying semiring, then introduce the polarity through two
further axioms, isolate the role of each, exhibit canonical models, and
verify that the resulting eleven-axiom system is independent.

\begin{definition}[Primitive semiring]\label{def:semiring}
A \emph{primitive semiring} is a structure $(X,\vee,\otimes,\ev,\eo)$, with
$\vee,\otimes$ binary operations and $\ev,\eo \in X$, satisfying for all
$a,b,c \in X$:
\begin{align}
  a \vee a &= a, \tag{J0}\label{ax:J0}\\
  (a \vee b) \vee c &= a \vee (b \vee c), \tag{J1}\label{ax:J1}\\
  a \vee b &= b \vee a, \tag{J2}\label{ax:J2}\\
  a \vee \ev &= a, \tag{J3}\label{ax:J3}\\
  (a \otimes b) \otimes c &= a \otimes (b \otimes c), \tag{T1}\label{ax:T1}\\
  a \otimes b &= b \otimes a, \tag{T2}\label{ax:T2}\\
  a \otimes \eo &= a, \tag{T3}\label{ax:T3}\\
  a \otimes (b \vee c) &= (a \otimes b) \vee (a \otimes c),
    \tag{TJ1}\label{ax:TJ1}\\
  a \otimes \ev &= \ev. \tag{TJ2}\label{ax:TJ2}
\end{align}
\end{definition}

Axioms \eqref{ax:J0} to \eqref{ax:J3} make $(X,\vee,\ev)$ an idempotent
commutative monoid, axioms \eqref{ax:T1} to \eqref{ax:T3} make
$(X,\otimes,\eo)$ a commutative monoid, and \eqref{ax:TJ1} to \eqref{ax:TJ2}
link the two, so that $(X,\vee,\otimes,\ev,\eo)$ is an idempotent commutative
semiring \citep{Golan1999}. These nine axioms describe the algebraic layer
alone; no duality, reflection, or residual is present.

\begin{definition}[Polar Semiring]\label{def:ps}
A \emph{Polar Semiring} is a structure $(X,\vee,\otimes,{}^{*},\ev,\eo)$
whose reduct $(X,\vee,\otimes,\ev,\eo)$ is a primitive semiring and whose
unary operation ${}^{*}\colon X \to X$ satisfies, for all $a,b \in X$:
\begin{align}
  \dstar{(\dstar{a})} &= a, \tag{P}\label{ax:P}\\
  a \vee b = b \;&\Longleftrightarrow\; \dstar{b} \vee \dstar{a} = \dstar{a}.
    \tag{PJ}\label{ax:PJ}
\end{align}
The operation ${}^{*}$ is the \emph{polarity} of the structure. We
abbreviate ``Polar Semiring'' by PS.
\end{definition}

% Axioms \eqref{ax:P} and \eqref{ax:PJ} are of a different nature from the
% nine semiring axioms: they add no algebraic operation but introduce a
% primitive polarity and its single point of contact with the semiring. To
% make that contact explicit, define the \emph{join order} by
% \begin{equation}\label{eq:orders}
%   a \leq b \;:\Longleftrightarrow\; a \vee b = b,
% \end{equation}
% which is a partial order by \eqref{ax:J0} to \eqref{ax:J2}, with least
% element $\ev$ by \eqref{ax:J3} and $a \vee b$ the join of $a$ and $b$. Under
% \eqref{eq:orders}, axiom \eqref{ax:PJ} states that $a \leq b$ if and only if
% $\dstar{b} \leq \dstar{a}$; that is, ${}^{*}$ is an antitone involution on
% $(X,\leq)$, the classical notion of a polarity on an ordered set
% \citep{Birkhoff1940}. No dual operation belongs to the signature, and none
% is constructed here; the structures generated from ${}^{*}$ by conjugation
% are the subject of \Cref{sec:arch}. The next two propositions record the
% distinct roles of \eqref{ax:P} and \eqref{ax:PJ}.



\begin{proposition}[Role of the polarity axiom]\label{prop:role-P}
Under \eqref{ax:P}, the polarity ${}^{*}$ is a bijection with
${}^{*} = {}^{*}{}^{-1}$.
\end{proposition}

\begin{proof}
By \eqref{ax:P} each $a$ equals $\dstar{(\dstar{a})}$, so ${}^{*}$ is
surjective; and $\dstar{a} = \dstar{b}$ gives
$a = \dstar{(\dstar{a})} = \dstar{(\dstar{b})} = b$, so ${}^{*}$ is
injective. The same equation shows ${}^{*}$ equals its own inverse.
\end{proof}

\begin{proposition}[Nontriviality of polarity]\label{prop:role-PJ}
Let $(X,\vee,{}^{*})$ satisfy \eqref{ax:J0} to \eqref{ax:J2} and
\eqref{ax:PJ}. If ${}^{*} = \mathrm{id}_X$, then $X$ is a singleton.
Consequently, on every nontrivial Polar Semiring, ${}^{*} \neq \mathrm{id}_X$.
\end{proposition}

\begin{proof}
Assume ${}^{*} = \mathrm{id}_X$, let $a,b \in X$ be arbitrary, and set
$c = a \vee b$. By \eqref{ax:J1} and \eqref{ax:J0},
$$
  a \vee c = a \vee (a \vee b) = (a \vee a) \vee b = a \vee b = c.
$$
Applying \eqref{ax:PJ} to the pair $(a,c)$, the equality $a \vee c = c$
yields $\dstar{c} \vee \dstar{a} = \dstar{a}$, which under
${}^{*} = \mathrm{id}_X$ reads $c \vee a = a$; but \eqref{ax:J2} gives
$c \vee a = a \vee c = c$, so $a = c$. Symmetrically, by \eqref{ax:J1},
\eqref{ax:J0}, and \eqref{ax:J2},
$$
  b \vee c = b \vee (a \vee b) = (b \vee b) \vee a = b \vee a = a \vee b = c,
$$
and \eqref{ax:PJ} applied to $(b,c)$ gives $\dstar{c} \vee \dstar{b} =
\dstar{b}$, that is, $c \vee b = b$; again \eqref{ax:J2} gives $c \vee b =
b \vee c = c$, so $b = c$. Hence $a = c = b$. As $a,b$ were arbitrary, any
two elements of $X$ coincide, so $X$ is a singleton.
\end{proof}

The proof uses only the join idempotence, associativity, and commutativity
\eqref{ax:J0} to \eqref{ax:J2} together with \eqref{ax:PJ}; neither the join
unit \eqref{ax:J3}, nor any tensor axiom, nor the order or the derived meet
is invoked. Thus \eqref{ax:PJ} does more than relate the primitive semiring
to the involution: on every nontrivial model it forbids the trivial
involution, so that the polarity of a nontrivial Polar Semiring is a genuine
reflection rather than the identity. Axiom \eqref{ax:P} alone would permit
${}^{*} = \mathrm{id}_X$; it is \eqref{ax:PJ} that rules this out. The deeper
structural role of \eqref{ax:PJ}, generating the compatibility between a
Polar Semiring and its dual, is established in \Cref{sec:arch}.

We now exhibit canonical models, confirming that the class is nonempty and
that primitive polarity coexists with the semiring operations. In each case
the join order coincides with the natural order of a chain, so
\eqref{ax:J0} to \eqref{ax:J3} are immediate; we verify the remaining
axioms directly.

\begin{example}[Involutive semiring model]\label{ex:involutive-semiring}
An \emph{involutive semiring} consists of a structure
$(X,\vee,\otimes,{}^{*},e_{\vee},e_{\otimes})$,
where $(X,\vee,e_{\vee})$ and $(X,\otimes,e_{\otimes})$ are commutative monoids,
${}^{*}:X\to X$ is an involution satisfying
$(a^{*})^{*}=a$,
and distributes over both operations,
\[
(a\vee b)^{*}=a^{*}\vee b^{*},\qquad
(a\otimes b)^{*}=a^{*}\otimes b^{*}.
\]
Moreover, $\otimes$ distributes over $\vee$.

Compared with Polar Semirings, this structure does not require
$\vee$ to be idempotent and does not impose any compatibility between
${}^{*}$ and the order induced by $\vee$. In particular, the identity map
${}^{*}=\mathrm{id}_{X}$ always satisfies the involution and distribution
axioms. Hence involutive semirings illustrate that the semiring axioms
together with involution alone do not exclude the trivial involution,
highlighting the essential role of the polarity axiom \eqref{ax:PJ}.
\end{example}

\begin{example}[Boolean model]\label{ex:boolean}
Let $X = \{0,1\}$ with $\vee = \max$, $\otimes = \min$, $\dstar{x} = 1-x$,
$\ev = 0$, $\eo = 1$. The monoid axioms \eqref{ax:T1} to \eqref{ax:T3} hold
for $\min$ with unit $1$; \eqref{ax:TJ1} is the distributivity of $\min$
over $\max$; \eqref{ax:TJ2} reads $\min(a,0)=0$; \eqref{ax:P} holds since
$1-(1-x)=x$; and \eqref{ax:PJ} holds because $x \mapsto 1-x$ is antitone.
This is the smallest nontrivial Polar Semiring.
\end{example}

\begin{example}[Tropical model]\label{ex:tropical}
Let $X = \mathbb{R} \cup \{-\infty,+\infty\}$ with $\vee = \max$,
$a \otimes b = a+b$, $\dstar{x} = -x$, $\ev = -\infty$, $\eo = 0$, under the
convention $(-\infty)+x = -\infty$ for all $x$, including $x = +\infty$. Then
$(X,\otimes,0)$ is a commutative monoid, so \eqref{ax:T1} to \eqref{ax:T3}
hold; addition is monotone, so $a + \max(b,c) = \max(a+b,a+c)$ gives
\eqref{ax:TJ1}; the stated convention gives $a \otimes (-\infty) = -\infty$,
which is \eqref{ax:TJ2}; \eqref{ax:P} holds since $-(-x)=x$; and negation is
antitone, giving \eqref{ax:PJ}.
\end{example}

\begin{example}[Quadratic model on the unit interval]\label{ex:quadratic}
Let $X = [0,1]$ with $\vee = \max$, $a \otimes b = ab$, and
\begin{equation}\label{eq:quadratic-star}
  \dstar{x} := \bigl(1 - \sqrt{x}\bigr)^{2},
\end{equation}
with $\ev = 0$ and $\eo = 1$. Ordinary multiplication makes
$(X,\otimes,1)$ a commutative monoid and is monotone, giving \eqref{ax:T1}
to \eqref{ax:T3} and \eqref{ax:TJ1}; $a \otimes 0 = 0$ gives \eqref{ax:TJ2}.
The map \eqref{eq:quadratic-star} is strictly decreasing with $\dstar{0}=1$,
$\dstar{1}=0$, whence \eqref{ax:PJ}, and it is an involution: since
$1 - \sqrt{x} \geq 0$ on $[0,1]$,
$$
  \dstar{(\dstar{x})}
  = \Bigl(1 - \sqrt{(1-\sqrt{x})^{2}}\Bigr)^{2}
  = \bigl(1 - (1 - \sqrt{x})\bigr)^{2}
  = x,
$$
which is \eqref{ax:P}. This model plays a decisive role in
\Cref{subsec:impossibility}.
\end{example}

No dual semiring is introduced or analyzed at this stage; the three models
are given purely in the primitive signature $(X,\vee,\otimes,{}^{*},\ev,\eo)$.
We close by verifying that no axiom is redundant.

\begin{theorem}[Independence]\label{thm:independence}
The eleven-axiom system of \Cref{def:semiring} and \Cref{def:ps} is
independent: for each axiom there exists a finite structure
$(X,\vee,\otimes,{}^{*},\ev,\eo)$ satisfying the remaining ten axioms but not
the selected one.
\end{theorem}

\begin{proof}
For each axiom a finite countermodel was obtained by exhaustive search with
the model finder Mace4, and the nonderivability claims were checked
independently with the theorem prover Prover9 \citep{McCune2005}. The
countermodels have at most \noteAI{fill in the maximal cardinality} elements
and are listed in \noteAI{decide: appendix with the model tables, or
supplementary material}.
\noteAI{Per final rules \S6.2, the original Mace4/Prover9 runs used the
obsolete axiom names (M0 to M3, S1 to S3, SM1 to SM2, PM) and the obsolete
choice of primitives ($\wedge$ primitive). The runs must be repeated in the
present signature (primitives $\vee,\otimes,{}^{*},\ev,\eo$; axioms J0 to J3,
T1 to T3, TJ1, TJ2, P, PJ) before this theorem is submitted.}
\end{proof}

Each axiom therefore contributes an essential and independent ingredient to
the definition of a Polar Semiring. This section has shown that primitive
polarity may be adjoined to a semiring consistently and independently, yet
so far it is only an involution constrained by a single compatibility
axiom. Whether these minimal assumptions already determine a dual
architecture is the question of the next section, which answers it in the
affirmative: primitive polarity canonically generates a dual semiring
together with a compatible order.

% % =========================================================
% doc3-structures.tex — Section 3: Architectural Analysis
% Source: sec2-outline.md (content 1.0–1.2.3, MAPPED into the
% three-layer scheme of d0, per final-rules §1/§6.5) +
% d2-from-axioms-to-theorems-bo-cuc.md (Prop 3.1–3.4)
% Layer mapping: 3.1 Algebraic  ⟵ J0–J3, T1–T3, TJ1–TJ2
%                3.2 Reflection ⟵ P (free duality, both modules)
%                3.3 Integration⟵ PJ (order compatibility)
% =========================================================

\section{Duality on Polar Semiring}\label{sec:arch}

The purpose of this section is to explain why the eleven axioms of
\Cref{def:ps} decompose into a modular architecture, and to prove the
Fundamental Structure Theorem announced in
\Cref{thm:fundamental-announce}. The decomposition has three layers. The
\emph{algebraic layer} (\Cref{subsec:algebraic}) is the polarity-free
reduct: two monoid modules and the two axioms linking them. The
\emph{reflection layer} (\Cref{subsec:reflection}) shows that the involution
axiom \eqref{ax:P} alone reflects each module into a dual module and yields
every De Morgan law at no further axiomatic cost. The \emph{integration
layer} (\Cref{subsec:integration}) shows that the single axiom
\eqref{ax:PJ} integrates the reflected order with the original one,
producing a lattice and completing the structure. Throughout, recall from
\Cref{rem:three-content-axioms} that only \eqref{ax:PJ},
\eqref{ax:TJ1}, and \eqref{ax:TJ2} carry compatibility content; the
architecture below makes this precise.






% Two scale parameters: physical size per grid unit.
\def\xunit{0.8cm}
\def\yunit{0.5cm}

\begin{figure}[ht]
\centering
\begin{tikzpicture}[x=\xunit, y=\yunit, thick, font=\small]

% ===================================================================
% Integer grid. Each box given by bottom-left (BL) and top-right (TR).
% Outer boxes are 16 units wide; inner boxes are 6 units wide with a
% 2-unit gap, so BL/TR and all centers land on integers.
% Layout, bottom to top: Primal semiring / Polarity / Dual semiring.
% ===================================================================

% ---- Primal semiring: outer box (0,0) to (16,4) ----
\draw[dashed] (0,0) rectangle (16,4);
\node[anchor=north west, inner sep=3pt] at (0,4) {\itshape Primal semiring};

% Join: BL (1,1) TR (7,3), center (4,2)
\draw (1,1) rectangle (7,3);
\node[align=center] at (4,2) {Join (J0--J3)};

% Tensor: BL (9,1) TR (15,3), center (12,2)
\draw (9,1) rectangle (15,3);
\node[align=center] at (12,2) {Tensor (T1--T3)};

% intra-layer edge: Join -- Tensor, at mid-height y=2
\draw (7,2) -- (9,2);
\node[fill=white, inner sep=1pt] at (8,2) {\scriptsize (TJ1, TJ2)};

% ---- Polarity: BL (2,6) TR (14,8), center (8,7) ----
\draw (2,6) rectangle (14,8);
\node at (8,7) {Polarity $(\mathrm{P})$};

% ---- Dual semiring: outer box (0,10) to (16,14) ----
\draw[dashed] (0,10) rectangle (16,14);
\node[anchor=north west, inner sep=3pt] at (0,14) {\itshape Dual semiring};

% Meet: BL (1,11) TR (7,13), center (4,12)
\draw (1,11) rectangle (7,13);
\node at (4,12) {Meet};

% Plus: BL (9,11) TR (15,13), center (12,12)
\draw (9,11) rectangle (15,13);
\node at (12,12) {Plus};

% intra-layer edge: Meet -- Plus, at mid-height y=12
\draw (7,12) -- (9,12);
% \node[fill=white, inner sep=1pt] at (8,12) {\scriptsize (RS3)};

% ---- vertical links: outer box to outer box, both centered at x=8 ----
\draw[->] (8,4) -- (8,6) node[midway, fill=white, inner sep=1pt] {\scriptsize (PJ)};
% \draw[->] (8,4) -- (8,6);
\draw[->] (8,8) -- (8,10);

\end{tikzpicture}
\caption{The axiomatic modules of a Polar Semiring. The primal semiring
$(X,\vee,\otimes,e_\vee,e_\otimes)$ and its dual $(X,\wedge,\oplus,e_\wedge,e_\oplus)$
are related through the polarity, which supplies axioms (P) and (PJ).}
\label{fig:architecture}
\end{figure}


\subsection{Algebraic layer}\label{subsec:algebraic}

The algebraic layer consists of two independent modules and their
interaction:
\begin{enumerate}
  \item the \emph{join module} $(X,\vee,\ev)$, an idempotent commutative
  monoid by \eqref{ax:J0}--\eqref{ax:J3};
  \item the \emph{tensor module} $(X,\otimes,\eo)$, a commutative monoid by
  \eqref{ax:T1}--\eqref{ax:T3};
  \item the interaction axioms \eqref{ax:TJ1}--\eqref{ax:TJ2}, which make
  $(X,\vee,\otimes,\ev,\eo)$ an idempotent commutative semiring
  \citep{Golan1999}.
\end{enumerate}
The join module already carries order-theoretic information: by
\eqref{ax:J0}--\eqref{ax:J2}, the relation $\leq$ of \eqref{eq:orders} is a
partial order (reflexivity is \eqref{ax:J0}, antisymmetry follows from
\eqref{ax:J2}, transitivity from \eqref{ax:J1}), with least element $\ev$ by
\eqref{ax:J3}, and $a \vee b$ is the least upper bound of $a$ and $b$. The
interaction axioms give monotonicity of the tensor:

\begin{lemma}\label{lem:tensor-monotone}
In any Polar Semiring, $a \leq b$ implies $a \otimes c \leq b \otimes c$.
\end{lemma}

\begin{proof}
If $a \vee b = b$, then by \eqref{ax:TJ1} and \eqref{ax:T2},
$(a \otimes c) \vee (b \otimes c) = (a \vee b) \otimes c = b \otimes c$.
\end{proof}

Nothing in this layer mentions ${}^{*}$: it is ordinary semiring theory.

\subsection{Reflection layer}\label{subsec:reflection}

We now add the involution axiom \eqref{ax:P} and nothing else, and show
that it \emph{reflects} each module of the algebraic layer into a dual
module, with all duality laws obtained for free. We first fix terminology.
Following the historical order of the notions---involutions belong to pure
algebra, polarities to order theory \citep{Birkhoff1940}---we call a
structure \emph{involutive} when it satisfies only
$\dstar{(\dstar{a})} = a$, and reserve the adjective \emph{polar} for
structures in which the involution is additionally compatible with the
induced order via \eqref{ax:PJ}. Thus an \emph{involutive commutative
monoid} is a commutative monoid equipped with an involution; no order is
involved, and indeed at this level no order need exist.

\begin{proposition}[Duality laws]\label{prop:duality}
Let $(X,\vee,\ev)$ and $(X,\otimes,\eo)$ be commutative monoids and let
${}^{*}$ satisfy \eqref{ax:P}. With the derived operations
\eqref{eq:derived-ops}, the following hold for all $a,b \in X$:
\begin{align}
  \dstar{(a \wedge b)} &= \dstar{a} \vee \dstar{b},
    \label{eq:RD1}\\
  \dstar{(a \oplus b)} &= \dstar{a} \otimes \dstar{b},
    \label{eq:RD2}\\
  \dstar{(a \vee b)} &= \dstar{a} \wedge \dstar{b},
    \label{eq:RD3}\\
  \dstar{(a \otimes b)} &= \dstar{a} \oplus \dstar{b},
    \label{eq:RD4}\\
  \dstar{\ev} = \ew,\quad \dstar{\ew} = \ev,&\quad
  \dstar{\eo} = \ep,\quad \dstar{\ep} = \eo.
    \label{eq:RD5}
\end{align}
\end{proposition}

\begin{proof}
For \eqref{eq:RD1}, apply \eqref{ax:P} to the definition:
$\dstar{(a \wedge b)}
 = \dstar{(\dstar{(\dstar{a} \vee \dstar{b})})}
 = \dstar{a} \vee \dstar{b}$.
For \eqref{eq:RD3}, compute
$\dstar{a} \wedge \dstar{b}
 = \dstar{(\dstar{(\dstar{a})} \vee \dstar{(\dstar{b})})}
 = \dstar{(a \vee b)}$ by \eqref{ax:P}. Equations \eqref{eq:RD2} and
\eqref{eq:RD4} are identical computations with $\otimes,\oplus$ in place of
$\vee,\wedge$. In \eqref{eq:RD5}, the first and third equations are the
definitions in \eqref{eq:derived-ops}, and the other two follow by
\eqref{ax:P}.
\end{proof}

\begin{proposition}[Reflected modules]\label{prop:reflected-modules}
Under the hypotheses of \Cref{prop:duality}:
\begin{enumerate}
  \item $(X,\wedge,\ew)$ is a commutative monoid, the \emph{meet module};
  \item $(X,\oplus,\ep)$ is a commutative monoid, the \emph{plus module};
  \item if moreover \eqref{ax:J0} holds, then $\wedge$ is idempotent.
\end{enumerate}
\end{proposition}

\begin{proof}
Commutativity of $\wedge$ is immediate from \eqref{ax:J2}. For
associativity, using \eqref{eq:RD1},
\[
  (a \wedge b) \wedge c
  = \dstar{\bigl(\dstar{(a \wedge b)} \vee \dstar{c}\bigr)}
  = \dstar{\bigl((\dstar{a} \vee \dstar{b}) \vee \dstar{c}\bigr)},
\]
and \eqref{ax:J1} lets the inner join reassociate; reversing the steps gives
$a \wedge (b \wedge c)$. For the unit,
$a \wedge \ew
 = \dstar{(\dstar{a} \vee \dstar{\ew})}
 = \dstar{(\dstar{a} \vee \ev)}
 = \dstar{(\dstar{a})} = a$
by \eqref{eq:RD5}, \eqref{ax:J3}, and \eqref{ax:P}. Item (2) is the same
argument in the tensor module, using \eqref{ax:T1}--\eqref{ax:T3}. For (3),
$a \wedge a = \dstar{(\dstar{a} \vee \dstar{a})} = \dstar{(\dstar{a})} = a$
by \eqref{ax:J0} and \eqref{ax:P}.
\end{proof}

Two comments are in order. First, everything above is a consequence of the
involution equation alone; the reflection layer is \emph{axiomatically
free}. Second, at this stage the two relations $\leq$ and $\leqm$ of
\eqref{eq:orders} are both available, but nothing yet guarantees that they
coincide, nor that $\wedge$ computes greatest lower bounds with respect to
$\leq$. That is precisely the task of the third layer.

\subsection{Integration layer}\label{subsec:integration}

The final axiom \eqref{ax:PJ} is the first, and only, axiom relating the
polarity to the join module, and it is not derivable from the previous
layers (\Cref{thm:independence}). Its effect is to integrate the reflected
order with the original one.

\begin{proposition}[Order compatibility]\label{prop:order-compat}
In the presence of the other ten axioms, the following are equivalent:
\begin{enumerate}
  \item \eqref{ax:PJ}: $a \vee b = b \iff \dstar{b} \vee \dstar{a} = \dstar{a}$;
  \item $a \vee b = b \iff b \wedge a = a$;
  \item $a \leq b \iff \dstar{b} \leqm \dstar{a}$;
  \item the absorption laws:
  $a \wedge (a \vee b) = a$ and $a \vee (a \wedge b) = a$.
\end{enumerate}
\end{proposition}

\begin{proof}
(1)$\Leftrightarrow$(2): by \eqref{eq:derived-ops} and \eqref{ax:P},
$b \wedge a = a$ holds if and only if
$\dstar{(\dstar{b} \vee \dstar{a})} = a$, if and only if
$\dstar{b} \vee \dstar{a} = \dstar{a}$; so (2) is a literal restatement of
(1).

(1)$\Leftrightarrow$(3): unfolding \eqref{eq:orders},
$\dstar{b} \leqm \dstar{a}$ means $\dstar{b} \wedge \dstar{a} = \dstar{b}$,
that is, $\dstar{(b \vee a)} = \dstar{b}$ by \eqref{eq:RD3}, that is,
$b \vee a = a$ by injectivity of ${}^{*}$; renaming variables shows (3) to
be a restatement of (1) as well.

(2)$\Rightarrow$(4): since $a \leq a \vee b$ by
\eqref{ax:J0}--\eqref{ax:J1}, item (2) yields
$(a \vee b) \wedge a = a$, which is the first absorption law by
commutativity of $\wedge$. For the second, (2) applied in the direction
``$b \wedge a = a \Rightarrow a \vee b = b$'' with the pair
$(a \wedge b, a)$: indeed $a \wedge (a \wedge b) = a \wedge b$ by
idempotence and associativity of $\wedge$
(\Cref{prop:reflected-modules}), so $(a \wedge b) \vee a = a$, which is
the second absorption law.

(4)$\Rightarrow$(2): if $a \vee b = b$, then
$a \wedge b = a \wedge (a \vee b) = a$, so $b \wedge a = a$. Conversely, if
$b \wedge a = a$, then $a \vee b = (b \wedge a) \vee b = b$ by the second
absorption law and commutativity.
\end{proof}

\begin{proposition}[Order structure]\label{prop:order-structure}
In any Polar Semiring:
\begin{enumerate}
  \item $\leq$ is a partial order and $\leq \;=\; \leqm$;
  \item $\ev \leq a \leq \ew$ for all $a$;
  \item $a \leq b$ implies $a \wedge c \leq b \wedge c$ and
        $a \oplus c \leq b \oplus c$;
  \item $a \vee b$ is the least upper bound and $a \wedge b$ the greatest
        lower bound of $\{a,b\}$ with respect to $\leq$; consequently
        $(X,\vee,\wedge)$ is a (possibly non-distributive) lattice.
\end{enumerate}
\end{proposition}

\begin{proof}
(1) That $\leq$ is a partial order was noted in \Cref{subsec:algebraic}.
By \Cref{prop:order-compat}(2), $a \leq b$ if and only if $a \wedge b = a$,
which is $a \leqm b$; hence the two orders coincide.

(2) $\ev \leq a$ is \eqref{ax:J3} with \eqref{ax:J2}. For $a \leq \ew$: by
\eqref{ax:PJ} this is equivalent to
$\dstar{\ew} \vee \dstar{a} = \dstar{a}$, that is,
$\ev \vee \dstar{a} = \dstar{a}$, which holds by \eqref{ax:J3}.

(3) Suppose $a \leq b$, so $a \wedge b = a$ by
\Cref{prop:order-compat}(2). Then
$(a \wedge c) \wedge (b \wedge c) = (a \wedge b) \wedge (c \wedge c)
 = a \wedge c$
by \Cref{prop:reflected-modules}, whence $a \wedge c \leqm b \wedge c$, and
we conclude by (1). For $\oplus$: from $a \leq b$,
\eqref{ax:PJ} gives $\dstar{b} \leq \dstar{a}$, then
\Cref{lem:tensor-monotone} gives
$\dstar{b} \otimes \dstar{c} \leq \dstar{a} \otimes \dstar{c}$, and applying
\eqref{ax:PJ} once more together with \eqref{eq:RD2} yields
$a \oplus c \leq b \oplus c$.

(4) The least-upper-bound property of $\vee$ is standard in idempotent
commutative monoids: $a,b \leq a \vee b$, and if $a,b \leq c$ then
$(a \vee b) \vee c = c$. For $\wedge$, dualize: by \eqref{ax:PJ} and
\eqref{eq:RD3}, lower bounds of $\{a,b\}$ correspond under ${}^{*}$ to upper
bounds of $\{\dstar{a},\dstar{b}\}$, and the correspondence sends
$a \wedge b$ to $\dstar{a} \vee \dstar{b}$; the greatest-lower-bound
property follows. The absorption laws of \Cref{prop:order-compat}(4) then
exhibit $(X,\vee,\wedge)$ as a lattice, bounded by (2). Distributivity of
$\vee$ over $\wedge$ is not asserted and indeed fails in general.
\noteAI{If you want a concrete non-distributive example here, add the
diamond $M_3$ with a suitable involution and check
\eqref{ax:TJ1}--\eqref{ax:TJ2} for a chosen tensor; otherwise the claim
``fails in general'' should be softened or given a citation.}
\end{proof}

The structure obtained from the join module together with \eqref{ax:P} and
\eqref{ax:PJ} alone---a bounded, possibly non-distributive lattice with an
antitone involution satisfying the De Morgan laws---is what we call a
\emph{Polar Lattice}; it coincides with the classical notion of a bounded
i-lattice \citep{Kalman1958}. We return to this identification in the
orientation table of \Cref{sec:positioning}.

\subsection{The Fundamental Structure Theorem}\label{subsec:fundamental}

Assembling the three layers yields the main theorem of this section.

\begin{theorem}[Fundamental Structure Theorem]\label{thm:fundamental}
Let $(X,\vee,\otimes,{}^{*},\ev,\eo)$ be a Polar Semiring. Then:
\begin{enumerate}
  \item the \emph{dual structure} $(X,\wedge,\oplus,{}^{*},\ew,\ep)$ is
  again a Polar Semiring;
  \item the polarity ${}^{*}\colon X \to X$ is an isomorphism from
  $(X,\vee,\otimes,{}^{*},\ev,\eo)$ onto its dual;
  \item the join order of the original structure and the meet order induced
  by the dual structure coincide, in the sense that
  $\leq \;=\; \leqm$.
\end{enumerate}
\end{theorem}

\begin{proof}
(1) We verify the eleven axioms for
$(X,\wedge,\oplus,{}^{*},\ew,\ep)$. Axioms J0--J3 for $\wedge$ and T1--T3
for $\oplus$ are \Cref{prop:reflected-modules}. Axiom P is unchanged. For
the dual of \eqref{ax:TJ1}, using \eqref{eq:RD1}--\eqref{eq:RD4} and
\eqref{ax:TJ1} in the original structure,
\[
  a \oplus (b \wedge c)
  = \dstar{\bigl(\dstar{a} \otimes (\dstar{b} \vee \dstar{c})\bigr)}
  = \dstar{\bigl((\dstar{a} \otimes \dstar{b}) \vee
                 (\dstar{a} \otimes \dstar{c})\bigr)}
  = (a \oplus b) \wedge (a \oplus c).
\]
For the dual of \eqref{ax:TJ2},
$a \oplus \ew
 = \dstar{(\dstar{a} \otimes \ev)}
 = \dstar{\ev} = \ew$
by \eqref{eq:RD5} and \eqref{ax:TJ2}. Finally, the dual of \eqref{ax:PJ}
asserts $a \wedge b = b \iff \dstar{b} \wedge \dstar{a} = \dstar{a}$; by
\Cref{prop:order-compat}(2) and \Cref{prop:order-structure}(1), both sides
translate to statements about $\leq$, namely $b \leq a$ on the left and
$\dstar{a} \leq \dstar{b}$ on the right, and these are equivalent by
\eqref{ax:PJ}.

(2) The map ${}^{*}$ is a bijection by \eqref{ax:P}, and
\eqref{eq:RD3}--\eqref{eq:RD5} say exactly that it sends $\vee$ to
$\wedge$, $\otimes$ to $\oplus$, $\ev$ to $\ew$, and $\eo$ to $\ep$; it
commutes with itself trivially. Hence it is an isomorphism onto the dual
structure.

(3) This is \Cref{prop:order-structure}(1).
\end{proof}

\begin{remark}
The theorem justifies the architectural reading of the axiom system: the
algebraic layer supplies the raw modules, the reflection layer produces the
dual modules for free from \eqref{ax:P}, and the single integration axiom
\eqref{ax:PJ} welds the two orders together, at which point the duality
becomes a self-duality of the whole structure. In particular, applying the
construction twice returns the original Polar Semiring, since
$\dstar{(\dstar{a})} = a$.
\end{remark}

% =========================================================
% doc3-duality.tex — Section 3: The Duality Induced by Polarity
% Narrative: motivation / reflection / generation / integration /
% fundamental theorem / interpretation. Technical content from
% doc3-structures.tex, reorganized into a single argumentative flow.
% =========================================================

\section{The Duality Induced by Polarity}\label{sec:arch}

Section~\ref{sec:ps} established that a primitive polarity exists, is
logically independent of the semiring axioms, and is nontrivial on every
nontrivial model. It did not, however, reveal what the polarity produces.
The guiding question of the present section is therefore the following: does
a primitive polarity merely define an involution, or does it determine a
richer algebraic structure? We answer it in the affirmative, and the answer
unfolds in three consecutive steps. First the polarity reflects each
primitive operation into a dual operation; then the reflected operations are
shown to form a genuine semiring; finally the compatibility axiom
\eqref{ax:PJ} forces the primal and dual semirings to share a single order.
The hero of the section is thus not the polarity itself but the duality it
generates. Figure~\ref{fig:architecture} records the modules involved and
the axioms that link them.

% Two scale parameters: physical size per grid unit.
\def\xunit{0.8cm}
\def\yunit{0.5cm}

\begin{figure}[ht]
\centering
\begin{tikzpicture}[x=\xunit, y=\yunit, thick, font=\small]

% ---- Primal semiring: outer box (0,0) to (16,4) ----
\draw[dashed] (0,0) rectangle (16,4);
\node[anchor=north west, inner sep=3pt] at (0,4) {\itshape Primal semiring};
\draw (1,1) rectangle (7,3);
\node[align=center] at (4,2) {Join (J0--J3)};
\draw (9,1) rectangle (15,3);
\node[align=center] at (12,2) {Tensor (T1--T3)};
\draw (7,2) -- (9,2);
\node[fill=white, inner sep=1pt] at (8,2) {\scriptsize (TJ1, TJ2)};

% ---- Polarity: BL (2,6) TR (14,8), center (8,7) ----
\draw (2,6) rectangle (14,8);
\node at (8,7) {Polarity $(\mathrm{P})$};

% ---- Dual semiring: outer box (0,10) to (16,14) ----
\draw[dashed] (0,10) rectangle (16,14);
\node[anchor=north west, inner sep=3pt] at (0,14) {\itshape Dual semiring};
\draw (1,11) rectangle (7,13);
\node at (4,12) {Meet};
\draw (9,11) rectangle (15,13);
\node at (12,12) {Plus};
\draw (7,12) -- (9,12);

% ---- vertical links ----
\draw[->] (8,4) -- (8,6) node[midway, fill=white, inner sep=1pt] {\scriptsize (PJ)};
\draw[->] (8,8) -- (8,10);

\end{tikzpicture}
\caption{The modules of a PS. The primal semiring
$(X,\vee,\otimes,\ev,\eo)$ and its dual $(X,\wedge,\oplus,\ew,\ep)$ are
related through the polarity, which supplies axioms \eqref{ax:P} and
\eqref{ax:PJ}.}
\label{fig:architecture}
\end{figure}




\textcolor{blue}{we should reorganize:\\
Part 1 semiring --> idem com monoidal sinh order\\
Part2 add involution --> dual idem com semiring --> weak antitone\\
Part3 add PJ ---> merging 2 orders --> generate lattice.\\
Main Theorem: polar semiring generates polar semiring dual and antitone.}




\textcolor{blue}{Part 2 add involution}
The first step is reflection. We introduce the operations conjugate to the
primitive ones and record their basic laws; nothing beyond the involution
axiom \eqref{ax:P} is used.

\begin{definition}[Derived operations and meet order]\label{def:derived}
Let $(X,\vee,\otimes,{}^{*},\ev,\eo)$ be a PS. Define the \emph{meet}, the
\emph{plus}, and the derived units by
\begin{equation}\label{eq:derived-ops}
  a \wedge b := \dstar{(\dstar{a} \vee \dstar{b})},\qquad
  a \oplus b := \dstar{(\dstar{a} \otimes \dstar{b})},\qquad
  \ew := \dstar{\ev},\qquad
  \ep := \dstar{\eo},
\end{equation}
and the \emph{meet order} by
\begin{equation}\label{eq:meet-order}
  a \leqm b \;:\Longleftrightarrow\; a \wedge b = a.
\end{equation}
\end{definition}

A word on terminology fixes the two levels at which the involution operates.
Following the historical order of the two notions, with involutions
belonging to pure algebra and polarities to order theory
\citep{Birkhoff1940}, we call a structure \emph{involutive} when it satisfies
only $\dstar{(\dstar{a})} = a$, and reserve the adjective \emph{polar} for
structures in which the involution is additionally compatible with the
induced order via \eqref{ax:PJ}. The reflection step lives entirely at the
involutive level: no order is required for it. The duality laws below are the
De Morgan identities relating the primitive operations to their conjugates.

\begin{proposition}[Duality laws]\label{prop:duality}
Let $(X,\vee,\ev)$ and $(X,\otimes,\eo)$ be commutative monoids and let
${}^{*}$ satisfy \eqref{ax:P}. With the derived operations
\eqref{eq:derived-ops}, the following hold for all $a,b \in X$:
\begin{align}
  \dstar{(a \wedge b)} &= \dstar{a} \vee \dstar{b}, \label{eq:RD1}\\
  \dstar{(a \oplus b)} &= \dstar{a} \otimes \dstar{b}, \label{eq:RD2}\\
  \dstar{(a \vee b)} &= \dstar{a} \wedge \dstar{b}, \label{eq:RD3}\\
  \dstar{(a \otimes b)} &= \dstar{a} \oplus \dstar{b}, \label{eq:RD4}\\
  \dstar{\ev} = \ew,\quad \dstar{\ew} = \ev,&\quad
  \dstar{\eo} = \ep,\quad \dstar{\ep} = \eo. \label{eq:RD5}
\end{align}
\end{proposition}

\begin{proof}
For \eqref{eq:RD1}, apply \eqref{ax:P} to the definition:
$\dstar{(a \wedge b)}
 = \dstar{(\dstar{(\dstar{a} \vee \dstar{b})})}
 = \dstar{a} \vee \dstar{b}$.
For \eqref{eq:RD3},
$\dstar{a} \wedge \dstar{b}
 = \dstar{(\dstar{(\dstar{a})} \vee \dstar{(\dstar{b})})}
 = \dstar{(a \vee b)}$ by \eqref{ax:P}. Equations \eqref{eq:RD2} and
\eqref{eq:RD4} are the same computations with $\otimes,\oplus$ in place of
$\vee,\wedge$. In \eqref{eq:RD5}, the first and third equations are the
definitions in \eqref{eq:derived-ops}, and the other two follow by
\eqref{ax:P}.
\end{proof}

% \begin{proposition}[Reflected modules]\label{prop:reflected-modules}
% Under the hypotheses of \Cref{prop:duality}, $(X,\wedge,\ew)$ and
% $(X,\oplus,\ep)$ are commutative monoids, and if \eqref{ax:J0} holds then
% $\wedge$ is idempotent.
% \end{proposition}

% \begin{proof}

% \end{proof}

The polarity reflects every primitive object into a dual one, and
does so at no axiomatic cost beyond \eqref{ax:P}. This yields a candidate
dual semiring $(X,\wedge,\oplus,\ew,\ep)$, though at this stage only its two
monoid modules have been secured; the axioms linking them remain to be
verified. That is the second step, generation, and it too uses only
\eqref{ax:P}: the linking axioms of the dual are the reflections of the
linking axioms of the primal.

\begin{proposition}[Dual PS]\label{prop:dual-semiring}
Let $(X,\vee,\otimes,{}^{*},\ev,\eo)$ be a PS.
Then $(X,\wedge,\oplus,*, \ew,\ep)$
is also a PS.
\end{proposition}

\begin{proof}
Commutativity of $\wedge$ is immediate from \eqref{ax:J2}. For associativity,
using \eqref{eq:RD1},
$$
  (a \wedge b) \wedge c
  = \dstar{\bigl(\dstar{(a \wedge b)} \vee \dstar{c}\bigr)}
  = \dstar{\bigl((\dstar{a} \vee \dstar{b}) \vee \dstar{c}\bigr)},
$$
and \eqref{ax:J1} reassociates the inner join; reversing the steps gives
$a \wedge (b \wedge c)$. For the unit,
$a \wedge \ew
 = \dstar{(\dstar{a} \vee \dstar{\ew})}
 = \dstar{(\dstar{a} \vee \ev)}
 = \dstar{(\dstar{a})} = a$
by \eqref{eq:RD5}, \eqref{ax:J3}, and \eqref{ax:P}. The plus module is the
same argument with \eqref{ax:T1} to \eqref{ax:T3}. Idempotence is
$a \wedge a = \dstar{(\dstar{a} \vee \dstar{a})} = \dstar{(\dstar{a})} = a$
by \eqref{ax:J0} and \eqref{ax:P}.

Axioms \eqref{ax:J0} to \eqref{ax:J3} for $\wedge$ and \eqref{ax:T1} to
\eqref{ax:T3} for $\oplus$ are \Cref{prop:reflected-modules}. For the dual of
\eqref{ax:TJ1}, using \eqref{eq:RD1} to \eqref{eq:RD4} and \eqref{ax:TJ1} in
the primal,
$$
  a \oplus (b \wedge c)
  = \dstar{\bigl(\dstar{a} \otimes (\dstar{b} \vee \dstar{c})\bigr)}
  = \dstar{\bigl((\dstar{a} \otimes \dstar{b}) \vee
                 (\dstar{a} \otimes \dstar{c})\bigr)}
  = (a \oplus b) \wedge (a \oplus c).
$$
For the dual of \eqref{ax:TJ2},
$a \oplus \ew = \dstar{(\dstar{a} \otimes \ev)} = \dstar{\ev} = \ew$
by \eqref{eq:RD5} and \eqref{ax:TJ2}. Hence all nine semiring axioms hold for
the dual operations.
\end{proof}

Every PS thus canonically induces a dual semiring, the first principal result
of the section. The two semirings, however, have so far been produced side by
side, each with its own order, and it is natural to ask whether they are
independent. The third step, integration, shows that they are not: the single
axiom \eqref{ax:PJ}, the one axiom relating the polarity to the join and the
only one not yet used, welds the two together. We first record that the
tensor is monotone, which the integration will need on the dual side.



\textcolor{blue}{Part 3 order}

\textcolor{blue}{Part 1 semiring --> idem com monoidal sinh order}

First we shoulde notice that $(X, \vee, e_\vee)$ is an idempotent commutative monoidal, thus it induces an order $\leq$ defined by 
\begin{equation*}
    a \leq b \Leftrightarrow a \vee b = b.
\end{equation*}
This order clearly satisfies reflexivity, antisymmetry and transitivity according to \eqref{ax:J0}-\eqref{ax:J3}.

Since $(X, \wedge,1_\wedge)$ is an idempotent comutative monoidal, we may define 
\begin{equation*}
    a \leq^* b \Leftrightarrow a \wedge b = a.
\end{equation*}

By the definitions of these two orders, we have 
\begin{lemma}\label{lem:tensor-monotone}
For $a \leq b$ and $p\leq^* q$, we have
\begin{enumerate}
    \item $a \vee c \leq b \vee c$
    \item $a \otimes c \leq b \otimes c$
    \item $p \wedge c \leq^* q \wedge c$
    \item $p \oplus c \leq^* q \oplus c$
\end{enumerate}
\end{lemma}

\begin{proof}
If $a \vee b = b$, then by \eqref{ax:TJ1} and \eqref{ax:T2},
$(a \otimes c) \vee (b \otimes c) = (a \vee b) \otimes c = b \otimes c$.
\end{proof}

By definition of $\leq$ and $\leq^*$, we have 
$$a \leq b \iff \dstar{b} \leqm \dstar{a}.$$
It means that the antitone of polarity comes from the definition of orders we define.
In the following, we show that, the \eqref{ax:PJ} axiom is actualy equivalent to the fact that the two orders coincide
$$\leq \equiv \leq^*.$$
This makes the antitone of polarity happens on a single poset.

\begin{proposition}[Order compatibility]\label{prop:order-compat}
In the presence of the other ten axioms, the following are equivalent to \eqref{ax:PJ}
\begin{enumerate}
  \item $a \vee b = b \iff a \wedge b = a$;
  \item $a \leq b \iff \dstar{b} \leq \dstar{a}$;
  \item $a \wedge (a \vee b) = a$ and
        $a \vee (a \wedge b) = a$.
\end{enumerate}
The last equalities are known as the absorption laws of join and meet.
\end{proposition}

\begin{proof}
(1)$\Leftrightarrow$(2): by \eqref{eq:derived-ops} and \eqref{ax:P},
$b \wedge a = a$ holds if and only if
$\dstar{(\dstar{b} \vee \dstar{a})} = a$, if and only if
$\dstar{b} \vee \dstar{a} = \dstar{a}$; so (2) is a literal restatement of
(1).

(1)$\Leftrightarrow$(3): unfolding \eqref{eq:meet-order},
$\dstar{b} \leqm \dstar{a}$ means $\dstar{b} \wedge \dstar{a} = \dstar{b}$,
that is, $\dstar{(b \vee a)} = \dstar{b}$ by \eqref{eq:RD3}, that is,
$b \vee a = a$ by injectivity of ${}^{*}$; renaming the variables shows (3) to
be a restatement of (1).

(2)$\Rightarrow$(4): since $a \leq a \vee b$ by \eqref{ax:J0} and
\eqref{ax:J1}, item (2) yields $(a \vee b) \wedge a = a$, the first
absorption law. For the second, apply (2) in the direction
$b \wedge a = a \Rightarrow a \vee b = b$ to the pair $(a \wedge b, a)$: since
$a \wedge (a \wedge b) = a \wedge b$ by \Cref{prop:reflected-modules}, one has
$(a \wedge b) \vee a = a$.

(4)$\Rightarrow$(2): if $a \vee b = b$, then
$a \wedge b = a \wedge (a \vee b) = a$, so $b \wedge a = a$. Conversely, if
$b \wedge a = a$, then $a \vee b = (b \wedge a) \vee b = b$ by the second
absorption law and commutativity.
\end{proof}

\begin{proposition}[Order structure]\label{prop:order-structure}
In any PS:
\begin{enumerate}
%   \item $\leq$ is a partial order and $\leq \;=\; \leqm$;
  \item $\ev \leq a \leq \ew$ for all $a$;
%   \item $a \leq b$ implies $a \wedge c \leq b \wedge c$ and
%         $a \oplus c \leq b \oplus c$;
  \item $(X,\vee,\wedge)$ is a bounded, possibly non-distributive lattice.
\end{enumerate}
\end{proposition}

\begin{proof}
(1) That $\leq$ is a partial order was noted in Section~\ref{sec:ps}. By
\Cref{prop:order-compat}(2), $a \leq b$ if and only if $a \wedge b = a$, which
is $a \leqm b$; hence the two orders coincide.

(2) $\ev \leq a$ is \eqref{ax:J3} with \eqref{ax:J2}. For $a \leq \ew$:
by \eqref{ax:PJ} this is equivalent to
$\dstar{\ew} \vee \dstar{a} = \dstar{a}$, that is,
$\ev \vee \dstar{a} = \dstar{a}$, which holds by \eqref{ax:J3}.

(3) If $a \leq b$, then $a \wedge b = a$ by \Cref{prop:order-compat}(2), so
$(a \wedge c) \wedge (b \wedge c) = (a \wedge b) \wedge (c \wedge c)
 = a \wedge c$
by \Cref{prop:reflected-modules}, whence $a \wedge c \leqm b \wedge c$, and
(1) concludes. For $\oplus$: from $a \leq b$, \eqref{ax:PJ} gives
$\dstar{b} \leq \dstar{a}$, \Cref{lem:tensor-monotone} gives
$\dstar{b} \otimes \dstar{c} \leq \dstar{a} \otimes \dstar{c}$, and applying
\eqref{ax:PJ} once more with \eqref{eq:RD2} yields
$a \oplus c \leq b \oplus c$.

(4) The least-upper-bound property of $\vee$ is standard in idempotent
commutative monoids. For $\wedge$, dualize: by \eqref{ax:PJ} and
\eqref{eq:RD3}, the lower bounds of $\{a,b\}$ correspond under ${}^{*}$ to the
upper bounds of $\{\dstar{a},\dstar{b}\}$, and the correspondence sends
$a \wedge b$ to $\dstar{a} \vee \dstar{b}$; the greatest-lower-bound property
follows. The absorption laws of \Cref{prop:order-compat}(4) then exhibit
$(X,\vee,\wedge)$ as a lattice, bounded by (2). Distributivity of $\vee$ over
$\wedge$ is not asserted and indeed fails in general.
\noteAI{If you want a concrete non-distributive example here, add the diamond
$M_3$ with a suitable involution and check \eqref{ax:TJ1}--\eqref{ax:TJ2} for
a chosen tensor; otherwise the claim ``fails in general'' should be softened
or given a citation.}
\end{proof}

The join module together with \eqref{ax:P} and \eqref{ax:PJ} alone yields a
bounded, possibly non-distributive lattice carrying an antitone involution
that satisfies the De Morgan laws. We call such a structure a \emph{Polar
Lattice}; it coincides with the classical notion of a bounded i-lattice
\citep{Kalman1958}, an identification to which we return in
Section~\ref{sec:positioning}. The three steps now combine into the central
theorem of the section, which states that the induced duality is in fact a
self-duality.

\begin{theorem}[Fundamental Structure Theorem]\label{thm:fundamental}
Let $(X,\vee,\otimes,{}^{*},\ev,\eo)$ be a PS. Then the dual structure
$(X,\wedge,\oplus,{}^{*},\ew,\ep)$ is again a PS, the polarity
${}^{*}\colon X \to X$ is an isomorphism and antitone.
\end{theorem}

\begin{proof}
By \Cref{prop:dual-semiring} the dual reduct $(X,\wedge,\oplus,\ew,\ep)$ is a
primitive semiring, and axiom \eqref{ax:P} is unchanged under passage to the
dual. It remains to verify the dual of \eqref{ax:PJ}, namely
$a \wedge b = b \iff \dstar{b} \wedge \dstar{a} = \dstar{a}$; by
\Cref{prop:order-compat}(2) and \Cref{prop:order-structure}(1) both sides
translate into statements about $\leq$, namely $b \leq a$ on the left and
$\dstar{a} \leq \dstar{b}$ on the right, and these are equivalent by
\eqref{ax:PJ}. Thus the dual is a PS. The map ${}^{*}$ is a bijection by
\eqref{ax:P}, and \eqref{eq:RD3} to \eqref{eq:RD5} say exactly that it sends
$\vee$ to $\wedge$, $\otimes$ to $\oplus$, $\ev$ to $\ew$, and $\eo$ to
$\ep$; hence it is an isomorphism onto the dual. The coincidence of the two
orders is \Cref{prop:order-structure}(1).
\end{proof}

Section~\ref{sec:ps} proved that polarity can be taken as primitive; the
present section proves that, once primitive, polarity automatically generates
a complete dual architecture. Reflection produces the conjugate operations,
generation assembles them into a dual semiring, and integration fuses the two
semirings into a single self-dual structure, at which point applying the
construction twice returns the original PS, since $\dstar{(\dstar{a})} = a$.
What Section~\ref{sec:ps} presented as an isolated involution has become an
order-reversing isomorphism between a semiring and its dual, obtained from
three compatibility axioms and no residual. The natural question that remains
is where this generated duality sits among existing algebraic structures,
which is the subject of Section~\ref{sec:positioning}.



% =========================================================
% doc4-compare.tex — Section 4: Positioning
% Sole content source: d3-compare-PS-and-BGA.md
%   Appendix A → 4.1 ; Sec 2 → 4.2 (compressed: 1 theorem + one
%   11-row table, axiom names TRANSLATED to J/T/TJ/P/PJ) ;
%   Sec 3 → 4.3 (impossibility on [0,1]).
% The residual \res is the ONLY BGA-native symbol used, and only here.
% =========================================================

\section{Positioning within Existing Algebraic Frameworks}\label{sec:positioning}

The previous section showed that a primitive polarity generates a complete dual algebraic architecture. Since axioms \eqref{ax:P} and \eqref{ax:PJ} define an antitone involution, it is natural to ask whether this duality has already appeared in existing algebraic structures. Among the classical candidates, bounded Girard algebras are the closest to Polar Semirings. We show that every bounded Girard algebra admits a canonical Polar Semiring reduct, whereas the converse fails. Thus, primitive polarity is strictly more general than polarity derived from residuation.

% \begin{center}
% \small
% \begin{tabular}{@{}p{4.6cm}p{4.3cm}p{4.3cm}@{}}
% \toprule
% Class & Relation to PS & Axiomatic difference \\
% \midrule
% Polar Lattice ($=$ bounded i-lattice \citep{Kalman1958}) &
%   order reduct of every PS & none (reduct) \\
% De Morgan algebra \citep{Rasiowa1974} &
%   special case of the order reduct & adds $\vee/\wedge$ distributivity \\
% Boolean algebra &
%   special case of De Morgan algebras & adds excluded middle \\
% Semiring with involution \citep{Golan1999} &
%   PS minus \eqref{ax:PJ} & drops order compatibility \\
% Bounded Girard algebra &
%   proper subclass of PS (\Cref{cor:strict}) & adds a residual \\
% \bottomrule
% \end{tabular}
% \end{center}
\subsection{Polar semiring and bounded Girard algebra}
\label{subsec:bga-defn}
% \subsection{Polarity versus Residual}

We recall the definition of a bounded Girard algebra, rewritten entirely in
the notation of this paper. The residual $\res$ is the one primitive of the
definition that has no counterpart in a Polar Semiring; it is the only
symbol foreign to our signature, and it appears in this section only.

\begin{definition}[Bounded Girard algebra]\label{def:bga}
A \emph{bounded Girard algebra} (BGA) is a structure
\[
  (X,\vee,\wedge,\otimes,\res,\eo,\ep,\ew)
\]
satisfying:
\begin{align}
  &(X,\vee,\wedge) \text{ is a lattice}, \tag{BGA1}\label{ax:BGA1}\\
  &(X,\otimes,\eo) \text{ is a commutative monoid}, \tag{BGA2}\label{ax:BGA2}\\
  &a \otimes b \leq c \;\Longleftrightarrow\; b \leq a \res c,
     \tag{BGA3}\label{ax:BGA3}\\
  &\dstar{(\dstar{a})} = a, \quad\text{where } \dstar{a} := a \res \ep,
     \tag{BGA4}\label{ax:BGA4}\\
  &a \leq \ew \text{ for all } a, \tag{BGA5}\label{ax:BGA5}
\end{align}
where $\leq$ denotes the lattice order, $a \leq b :\iff a \vee b = b$.
The remaining operations of the PS signature are derived:
$\ev := \dstar{\ew}$ and $a \oplus b := \dstar{(\dstar{a} \otimes
\dstar{b})}$.
\end{definition}

\noteAI{\Cref{def:bga} follows the axiomatization recorded in the project
notes as sourced from \citep{Agliano2025}; verify against the published
formulation (in particular whether boundedness is stated via \eqref{ax:BGA5}
or via a bottom element) before submission.}

By \eqref{ax:BGA3}--\eqref{ax:BGA4}, the polarity of a BGA is a derived
operation, manufactured from the residual and the distinguished element
$\ep$. The question of this section is whether that manufacturing step is
essential.

% \subsection{Every bounded Girard algebra induces a Polar
% Semiring}\label{subsec:bga-implies-ps}

We first isolate the consequences of residuation that will be needed; all
of them use only \eqref{ax:BGA1}--\eqref{ax:BGA3}.

\begin{lemma}\label{lem:res-toolkit}
In any structure satisfying
\eqref{ax:BGA1}--\eqref{ax:BGA3}:
\begin{enumerate}
  \item $a \leq b \iff \eo \leq a \res b$;
  \item $b \otimes (b \res c) \leq c$;
  \item $a \leq b$ implies $a \otimes c \leq b \otimes c$;
  \item $a \leq b$ implies $b \res z \leq a \res z$; in particular, with
        $z := \ep$, the map ${}^{*}$ is antitone.
\end{enumerate}
\end{lemma}

\begin{proof}
(1) Put $b := \eo$ in \eqref{ax:BGA3} and use $a \otimes \eo = a$.
(2) By \eqref{ax:BGA3}, $b \otimes (b \res c) \leq c$ is equivalent to
$b \res c \leq b \res c$.
(3) By commutativity, $b \leq c \res (b \otimes c)$ holds unconditionally
(it is \eqref{ax:BGA3} applied to $c \otimes b \leq b \otimes c$); from
$a \leq b$ and transitivity, $a \leq c \res (b \otimes c)$, and
\eqref{ax:BGA3} converts this back to $a \otimes c \leq b \otimes c$.
(4) By (3), $a \otimes (b \res z) \leq b \otimes (b \res z)$, and by (2)
the right-hand side is $\leq z$; then \eqref{ax:BGA3} gives
$b \res z \leq a \res z$.
\end{proof}

\begin{theorem}[BGA $\Rightarrow$ PS]\label{thm:bga-implies-ps}
Let $(X,\vee,\wedge,\otimes,\res,\eo,\ep,\ew)$ be a bounded Girard algebra,
and let ${}^{*}$, $\ev$ be the derived operations of \Cref{def:bga}. Then
$(X,\vee,\otimes,{}^{*},\ev,\eo)$ is a Polar Semiring, and its derived meet
and order coincide with the lattice meet and order of the BGA.
\end{theorem}

\begin{proof}
\Cref{tab:eleven} lists, for each of the eleven PS axioms, the BGA facts
from which it follows. Seven rows are immediate: \eqref{ax:J0},
\eqref{ax:J1}, \eqref{ax:J2} are contained in the lattice axiom
\eqref{ax:BGA1}; \eqref{ax:T1}, \eqref{ax:T2}, \eqref{ax:T3} are
\eqref{ax:BGA2}; and \eqref{ax:P} is \eqref{ax:BGA4}. The four remaining
rows are proved after the table.

\begin{table}[ht]
\centering
\caption{Derivation of the eleven PS axioms from the BGA axioms.}
\label{tab:eleven}
\begin{tabular}{@{}llll@{}}
\toprule
PS axiom & Statement & Derived from & Status \\
\midrule
\eqref{ax:J0} & $a \vee a = a$ & \eqref{ax:BGA1} & immediate \\
\eqref{ax:J1} & associativity of $\vee$ & \eqref{ax:BGA1} & immediate \\
\eqref{ax:J2} & commutativity of $\vee$ & \eqref{ax:BGA1} & immediate \\
\eqref{ax:J3} & $a \vee \ev = a$ & \eqref{ax:BGA4}, \eqref{ax:BGA5},
  \Cref{lem:res-toolkit}(4) & proof below \\
\eqref{ax:T1} & associativity of $\otimes$ & \eqref{ax:BGA2} & immediate \\
\eqref{ax:T2} & commutativity of $\otimes$ & \eqref{ax:BGA2} & immediate \\
\eqref{ax:T3} & $a \otimes \eo = a$ & \eqref{ax:BGA2} & immediate \\
\eqref{ax:TJ1} & $\otimes$ distributes over $\vee$ &
  \eqref{ax:BGA3}, \Cref{lem:res-toolkit}(1,3) & proof below \\
\eqref{ax:TJ2} & $a \otimes \ev = \ev$ &
  \eqref{ax:BGA3}, \eqref{ax:J3} & proof below \\
\eqref{ax:P} & $\dstar{(\dstar{a})} = a$ & \eqref{ax:BGA4} & immediate \\
\eqref{ax:PJ} & polarity & \eqref{ax:BGA4},
  \Cref{lem:res-toolkit}(4) & proof below \\
\bottomrule
\end{tabular}
\end{table}

\emph{\eqref{ax:J3}.} By \eqref{ax:BGA5}, $a \leq \ew$ for all $a$; since
${}^{*}$ is antitone (\Cref{lem:res-toolkit}(4)),
$\ev = \dstar{\ew} \leq \dstar{a}$ for all $a$. As ${}^{*}$ is surjective by
\eqref{ax:BGA4}, every element of $X$ is of the form $\dstar{a}$, so $\ev$
is the least element of the lattice, and $a \vee \ev = a$ follows.

\emph{\eqref{ax:TJ1}.} The inequality
$(a \otimes b) \vee (a \otimes c) \leq a \otimes (b \vee c)$ follows from
\Cref{lem:res-toolkit}(3) applied to $b, c \leq b \vee c$. Conversely, put
$u := (a \otimes b) \vee (a \otimes c)$. From $a \otimes b \leq u$ and
$a \otimes c \leq u$, axiom \eqref{ax:BGA3} yields
$b \leq a \res u$ and $c \leq a \res u$, hence
$b \vee c \leq a \res u$, and \eqref{ax:BGA3} converts this to
$a \otimes (b \vee c) \leq u$.

\emph{\eqref{ax:TJ2}.} Since $\ev$ is the least element,
$\ev \leq a \otimes \ev$; and $a \otimes \ev \leq \ev$ is, by
\eqref{ax:BGA3}, equivalent to $\ev \leq a \res \ev$, which holds because
$\ev$ is least. Antisymmetry concludes.

\emph{\eqref{ax:PJ}.} Unfolding \eqref{eq:orders}, the axiom asserts
$a \leq b \iff \dstar{b} \leq \dstar{a}$. The forward direction is
\Cref{lem:res-toolkit}(4). For the converse, apply
\Cref{lem:res-toolkit}(4) to $\dstar{b} \leq \dstar{a}$ to obtain
$\dstar{(\dstar{a})} \leq \dstar{(\dstar{b})}$, and use \eqref{ax:BGA4}.

Finally, the derived meet of the induced PS coincides with the lattice
meet: for the antitone bijection ${}^{*}$, the element
$\dstar{(\dstar{a} \vee \dstar{b})}$ is the greatest lower bound of
$\{a,b\}$ in $(X,\leq)$, which in the lattice \eqref{ax:BGA1} is
$a \wedge b$.
\end{proof}

% \subsection{The impossibility theorem}\label{subsec:impossibility}

We now show that \Cref{thm:bga-implies-ps} cannot be reversed: the quadratic
model of \Cref{ex:quadratic} is a Polar Semiring that admits no bounded
Girard algebra structure. In a BGA, the polarity is by definition
$\dstar{a} = a \res \ep$ for the residual $\res$ of $\otimes$; on
$([0,1],\leq)$ with $\otimes$ the ordinary product, the residual, when it
exists, is uniquely determined by \eqref{ax:BGA3}:
\begin{equation}\label{eq:interval-residual}
  a \res c \;=\; \max\{\, b \in [0,1] : ab \leq c \,\}
          \;=\; \min\bigl(1,\, c/a\bigr)
  \quad (a \in (0,1]),
  \qquad 0 \res c = 1.
\end{equation}
Thus extendability of the quadratic model to a BGA amounts to the existence
of a dualizing constant $q \in [0,1]$ with $\dstar{a} = a \res q$ for all
$a$, and the theorem rules this out.

\begin{theorem}[Impossibility]\label{thm:impossibility}
Let $([0,1],\vee,\otimes,{}^{*},0,1)$ be the quadratic model of
\Cref{ex:quadratic}. There exists no $q \in [0,1]$ such that
\begin{equation}\label{eq:dagger}
  a \otimes b \leq q
  \;\Longleftrightarrow\;
  b \leq \dstar{a}
  \qquad\text{for all } a,b \in [0,1].
\end{equation}
Consequently, the quadratic model cannot be extended to a bounded Girard
algebra: no choice of residual and dualizing element $\ep$ makes its
polarity \eqref{eq:quadratic-star} arise as $a \mapsto a \res \ep$.
\end{theorem}

\begin{proof}
Fix $a \in (0,1]$. As functions of $b$, both sides of \eqref{eq:dagger} are
threshold conditions: the left-hand side holds precisely for
$b \leq \min(1, q/a)$, the right-hand side precisely for
$b \leq \dstar{a}$. Since \eqref{eq:dagger} must hold for every
$b \in [0,1]$, the thresholds must agree:
\begin{equation}\label{eq:threshold}
  \min\bigl(1,\, q/a\bigr) = \dstar{a}
  \qquad\text{for all } a \in (0,1].
\end{equation}
Evaluating \eqref{eq:threshold} at $a = 1$ gives
$\min(1,q) = \dstar{1} = 0$; as $q \in [0,1]$, this forces $q = 0$. But
substituting $q = 0$ into \eqref{eq:threshold} at any $a \in (0,1)$ gives
$0 = \dstar{a} = (1-\sqrt{a})^{2} > 0$, a contradiction. Hence no
$q$ satisfies \eqref{eq:dagger}.

For the final claim: in any BGA structure on this carrier extending the
given $(\vee,\otimes)$, the order is the usual order of $[0,1]$, the
residual is forced to be \eqref{eq:interval-residual} by \eqref{ax:BGA3},
and the polarity is $a \res \ep$ for the primitive constant
$q := \ep$; the displayed equivalence \eqref{eq:dagger} is then exactly the
requirement that this derived polarity coincide with
\eqref{eq:quadratic-star}.
\noteAI{Per final-rules \S6.1, step \eqref{eq:threshold} (the
``$(\dagger)$'' step of the source notes) required independent
verification. I have re-derived it above and it checks out: for fixed
$a\in(0,1]$ the LHS of \eqref{eq:dagger} is $b\le\min(1,q/a)$, the RHS is
$b\le\dstar{a}$, equality of the two down-sets on $[0,1]$ forces equality
of thresholds; $a=1$ forces $q=0$; $q=0$ fails for every $a<1$. Please
perform your own final check (and, if desired, a Mace4 confirmation on a
finite subchain) before submission.}
\end{proof}

\begin{corollary}\label{cor:strict}
$\BGA \subsetneq \PS$: the class of Polar Semirings induced by bounded
Girard algebras via \Cref{thm:bga-implies-ps} is a proper subclass of the
class of all Polar Semirings.
\end{corollary}

\begin{remark}\label{rem:why-impossible}
The impossibility has a structural explanation. Among the eleven PS axioms,
none constrains $\otimes$ and ${}^{*}$ jointly through the order:
\eqref{ax:TJ1}--\eqref{ax:TJ2} relate $\otimes$ to $\vee$ without
mentioning ${}^{*}$, while \eqref{ax:P} and \eqref{ax:PJ} relate ${}^{*}$
to $\vee$ without mentioning $\otimes$. Whether a given polarity is
residuation-compatible with a given tensor is therefore left entirely
undetermined by the axioms, and \Cref{thm:impossibility} shows that this
freedom is genuine: the gap between the two classes is non-empty. In this
precise sense, Polar Semirings are a residuation-free generalization of
bounded Girard algebras---the polarity is promoted to a primitive, the
residual is discarded, and \Cref{thm:bga-implies-ps,thm:impossibility}
together certify that the generalization is strict rather than a
reformulation.
\end{remark}

\subsection{Polar semiring and other structures}
\label{subsec:related-structures}

The gap isolated in \Cref{thm:bga-implies-ps,thm:impossibility} is not
peculiar to bounded Girard algebras. It reflects a single structural fact
about the eleven axioms: \eqref{ax:PJ} ties the polarity to the order
without reference to $\otimes$, so any framework that manufactures the
polarity from $\otimes$ and an extra primitive --- a residual, a
dualizing constant, a fixed negation --- adds information that
\eqref{ax:J0}--\eqref{ax:TJ2} alone do not supply. The remaining
comparisons below trace the same fault line through three further
traditions, without repeating the reduct-and-impossibility argument in
each case.

The shape of \Cref{def:bga} is itself inherited from linear logic: a
bounded Girard algebra is an algebraic distillation of the phase
semantics of linear logic \noteAI{cite Girard's original phase-semantics
paper, e.g. Girard 1987}, and is order-equivalent to a
\emph{Girard quantale} --- a join-complete lattice-ordered monoid carrying
a dualizing element \noteAI{cite a Girard-quantale reference, e.g.
Rosenthal or rosenthal1990quantales}. Both traditions treat
the polarity as something derived from residuation on a complete lattice,
which is exactly the feature \Cref{subsec:bga-defn} isolates and rejects;
we do not unpack them separately here, since \eqref{ax:BGA1}--\eqref{ax:BGA5}
already exhibit their defining move in the sharpest available form.

A different departure is taken by involutive semirings, studied purely as
algebraic objects rather than as models of a logic
\citep{dolinka2000minimal,dheena2018ideals}. There the polarity is a
primitive anti-automorphism of $\otimes$, but it is required to interact
with $\otimes$ alone: nothing analogous to \eqref{ax:PJ} ties it to the
order, and $\vee$ is not assumed idempotent, so \eqref{ax:J0} fails as
well. The additive and multiplicative reducts collapse into a single
pair of operations, $\wedge \equiv \vee$ and $\oplus \equiv \otimes$,
leaving no room for the order-compatible duality that
\eqref{ax:P}--\eqref{ax:PJ} express. Polar Semirings and involutive
semirings are consequently incomparable rather than nested: each satisfies
axioms the other does not.

Dual pairs of semirings arise in the study of fuzzy sets and
residuated structures \citep{movckovr2025unification}, and sit on the
opposite side of PS from involutive semirings. All eleven axioms hold,
and the polarity remains primitive as in PS, but a dual pair of
semirings additionally assumes completeness of the order and demands that
$\otimes$ and $\vee$ distribute over a primitive $\wedge$, together with
distributivity of $\vee$ over $\wedge$. Where PS keeps $\wedge$ and
$\oplus$ derived and imposes no completeness, a dual pair of semirings
promotes them to primitives governed by their own axioms. It is thus a
proper strengthening of PS built on the same primitive polarity, in
contrast to bounded Girard algebras, which weaken the polarity to a
derived notion while adding the residual.

Table~\ref{tab:comparison} and Table~\ref{tab:axiom-comparison} summarize
these four comparisons side by side, recording for each structure which
operations are primitive versus derived and which of the eleven PS
axioms it satisfies.




\begin{table}[!ht]
\centering
\caption{Comparison of Polar Semirings, Bounded Girard Algebras, Dual Pairs of Semirings, and Classical Involutive Semirings}
\label{tab:comparison}
\renewcommand{\arraystretch}{1.2}
\begin{tabular}{lcccc}
\hline
\textbf{Property}
& \textbf{PS}
& \textbf{BGA}
& \textbf{DPS}
& \textbf{IS} \\
\hline

Residual
& ---
& Primitive
& ---
& --- \\

Polar $(*)$
& Primitive
& Derived
& Primitive
& Primitive \\

Join $(\vee)$
& Primitive
& Primitive
& Primitive
& Primitive \\

Tensor $(\otimes)$
& Primitive
& Primitive
& Primitive
& Primitive  \\

Meet $(\wedge)$
& Derived
& Primitive
& Primitive
& $\equiv \vee$ \\

Plus $(\oplus)$
& Derived
& Derived
& Primitive
& $\equiv \otimes$ \\

Unit of join $(1_\vee)$
& Primitive
& Primitive
& Primitive
& Primitive \\


Unit of tensor $(1_\otimes)$
& Primitive
& Primitive
& Primitive
& Primitive \\

Unit of meet $(1_\wedge)$
& Derived
& Primitive
& Primitive
& $=1_\vee$ \\

Unit of plus $(1_\oplus)$
& Derived
& Primitive
& Primitive
& $=1_\otimes$ \\

\hline

$\otimes$ distributes over $\vee$
& Primitive
& Derived
& Primitive
& Primitive \\

$\otimes$ distributes over $\wedge$
& ---
& ---
& Primitive
& Primitive \\

$\vee$ distributes over $\wedge$
& ---
& ---
& Primitive
& --- \\

involutive
& Primitive
& Primitive
& Primitive
& Primitive \\

antitone
& Derived
& Derived
& Derived
& --- \\

\hline

PS axioms satisfied
& All
& All
& All
& All but J0, PJ \\

% Additional assumptions
% & ---
% & ---
% & ---
% & $\wedge \equiv \vee$; $\oplus \equiv \otimes$ \\

\hline
\end{tabular}
\end{table}

% =========================================================
% doc5-summary.tex — Section 5: Conclusion
% New content (final-rules §1: Section 5, short).
% NOTE: d0-positioning-objectives.md still lacks an objective line for
% this section; suggested objective is stated in the first sentence and
% should be back-ported to d0 (final-rules §6.6).
% =========================================================

\section{Conclusion}\label{sec:conclusion}

We have introduced Polar Semirings, an axiomatization of polarity as a
primitive algebraic notion, and established their basic theory and their
position in the landscape of involutive structures. Three results summarize
the paper. First, the eleven axioms of \Cref{def:ps} are independent and
admit natural models on chains, including Boolean, tropical, and quadratic
models (\Cref{sec:ps}). Second, the axioms decompose into an algebraic
layer, a free reflection layer generated by the involution, and a one-axiom
integration layer, and this architecture yields the Fundamental Structure
Theorem: every Polar Semiring is canonically self-dual, with the polarity
acting as an isomorphism onto the dual structure and the two induced orders
coinciding (\Cref{thm:fundamental}). Third, every bounded Girard algebra
induces a Polar Semiring upon forgetting its residual
(\Cref{thm:bga-implies-ps}), whereas the quadratic model on the unit
interval admits no bounded Girard algebra extension
(\Cref{thm:impossibility}); hence $\BGA \subsetneq \PS$
(\Cref{cor:strict}), and the generalization is strict.

The conceptual message is that the dual architecture usually attributed to
residuation---dual monoids, De Morgan laws, order reversal---is in fact
generated by three compatibility axioms, \eqref{ax:PJ}, \eqref{ax:TJ1}, and
\eqref{ax:TJ2}, which are jointly much weaker than the residuation law
\eqref{ax:BGA3}.

Several directions remain open. On the axiomatic side, one may ask for the
exact strength of residuation over the PS axioms: which additional axioms,
formulated in the PS signature, characterize the Polar Semirings that do
extend to bounded Girard algebras? On the model-theoretic side, the variety
generated by Polar Semirings invites study of its free objects, its
subdirectly irreducible members, and the decidability of its equational
theory. On the structural side, the non-commutative case, in which
\eqref{ax:T2} is dropped and the polarity may split into left and right
polarities, and the distributive case, in which the induced lattice is
required to be distributive, both appear tractable with the methods of
\Cref{sec:arch}.
\noteAI{If you have partial results on any of these directions (e.g.\ the
characterization of residuable Polar Semirings), consider replacing the
corresponding sentence with a precise conjecture.}


\bibliographystyle{abbrvnat}
\bibliography{head-refs}

\end{document}
